\ifdefined\gkver\else\def\gkver{student}\fi
\documentclass[\gkver]{gaokaozhenti}

% 自感与涡流·高考真题专题练习（学生版、教师版共用此文件，编译方法与各卷相同）
% 题干、选项、答案和配图原样取自 YLXDXX/gaokao-physics-archive 提交 5a60bb3 的真题源文件；
% 每题前的注释写明出处和配图对应关系，相对上游的少量订正见仓库根目录 专题练习/README.md。

% 题目出处：（年份·试卷·原卷题号）
\newcommand{\source}[1]{{\sffamily（#1）}}

\begin{document}

\chapter{自感与涡流·高考真题专题练习}

\noindent 说明：本练习共 28 题，全部为高考真题原题，配图取自原卷；题号后括号内注明年份、试卷和原卷题号。

\section{自感现象}
\begin{enumerate}
% ---------------------------------------------------------------------------
% 第 1 题：2010 年海南卷第 7 题（上游 试卷/2010/海南/海南.tex 中 number 为 7 的一题）
% 配图：无
\item
%% number: 7
%% paperName: 2010年海南高考第 7 题 4 分
%% typeId: 2
%% type: 多选题
%% chapter: 电磁感应
%% point: 自感与互感，涡流
%% method: 无
%% score: 4
%% degree: 650
%% duplicateId: 0
%% body:
\source{2010·海南卷·7}下列说法正确的是\xzanswer{AC}

\fourchoices[answer=AC]
{当线圈中电流不变时，线圈中没有自感电动势}
{当线圈中电流反向时，线圈中自感电动势的方向与线圈中原电流的方向相反}
{当线圈中电流增大时，线圈中自感电动势的方向与线圈中电流的方向相反}
{当线圈中电流减小时，线圈中自感电动势的方向与线圈中电流的方向相反}

%% answer: AC
%% memo:
\memoanswer{
由法拉第电磁感应定律可知，当线圈中电流不变时，不产生自感电动势，A 对；

当线圈中电流反向时，相当于电流减小，线圈中自感电动势的方向与线圈中原电流的方向相同，B 错；

当线圈中电流增大时，自感电动势阻碍电流的增大，线圈中自感电动势的方向与线圈中电流的方向相反，C 对；

当线圈中电流减小时，自感电动势阻碍电流的减小，线圈中自感电动势的方向与线圈中电流的方向相同，D 错。
}

% ---------------------------------------------------------------------------
% 第 2 题：2007 年海南卷第 3 题（上游 试卷/2007/海南理综/海南理综.tex 中 number 为 3 的一题）
% 配图：02.svg←试卷/2007/海南理综/figs/03.svg
\item
%% number: 3
%% paperName: 2007年海南理综第 3 题 3 分
%% typeId: 1
%% type: 单选题
%% chapter: 电磁感应
%% point: 自感与互感，涡流
%% method: 无
%% score: 3
%% degree: 650
%% duplicateId: 0
%% body:
\source{2007·海南卷·3}在如图所示的电路中，a、b 为两个完全相同的灯泡，$L$ 为自感线圈，$E$ 为电源，S 为开关。关于两灯泡点亮和熄灭的先后次序，下列说法正确的是\xzanswer{C}

\onepicture[w=3.6cm]{\includesvg{figs/02}}

\fourchoices[answer=C]
{合上开关，a 先亮，b 后亮；断开开关，a、b 同时熄灭}
{合上开关，b 先亮，a 后亮；断开开关，a 先熄灭，b 后熄灭}
{合上开关，b 先亮，a 后亮；断开开关，a、b 同时熄灭}
{合上开关，a、b 同时亮；断开开关，b 先熄灭，a 后熄灭}

%% answer: C
%% memo:
\memoanswer{
由于自感作用，合上开关时，b 立即亮，而 a 所在支路开始时电流受到的阻碍很大，电流值很小，而后电流正常，故 a 后亮；断开开关，a、b 和 $L$ 组成一个闭合回路，电路中电流处处相等，稍后同时熄灭，故 C 正确。
}

% ---------------------------------------------------------------------------
% 第 3 题：2023 年北京卷第 5 题（上游 试卷/2023/北京/北京.tex 中 number 为 5 的一题）
% 配图：03.png←试卷/2023/北京/figs/05.png
\item
%% number: 5
%% paperName: 2023年北京高考第 5 题 3 分
%% typeId: 1
%% type: 单选题
%% chapter: 电磁感应
%% point: 自感与互感，涡流
%% method: 无
%% score: 3
%% degree: 750
%% duplicateId: 0
%% body:
\source{2023·北京卷·5}如图所示，$L$ 是自感系数很大、电阻很小的线圈，P、Q 是两个相同的小灯泡，开始时，开关 S 处于闭合状态，P 灯微亮，Q 灯正常发光，断开开关\xzanswer{D}
\onepicture[w=4.5cm]{figs/03.png}
\fourchoices[answer=D]
{P 与 Q 同时熄灭}
{P 比 Q 先熄灭}
{Q 闪亮后再熄灭}
{P 闪亮后再熄灭}
%% answer: D
%% memo:
\memoanswer{
由题知，开始时，开关 S 闭合时，由于 $L$ 的电阻很小，Q 灯正常发光，P 灯微亮，断开开关前通过 Q 灯的电流远大于通过 P 灯的电流，断开开关时，Q 所在电路未闭合，立即熄灭，由于自感，$L$ 中产生感应电动势，与 P 组成闭合回路，故 P 灯闪亮后再熄灭。

故选 D。
}

% ---------------------------------------------------------------------------
% 第 4 题：2017 年北京卷第 19 题（上游 试卷/2017/北京理综/北京理综.tex 中 number 为 19 的一题）
% 配图：04a.png←试卷/2017/北京理综/figs/19a.png；04b.png←试卷/2017/北京理综/figs/19b.png
\item
%% number: 19
%% paperName: 2017年北京理综第 19 题 6 分
%% typeId: 1
%% type: 单选题
%% chapter: 电磁感应
%% point: 自感与互感，涡流
%% method: 无
%% score: 6
%% degree: 650
%% duplicateId: 0
%% body:
\source{2017·北京卷·19}图 1 和图 2 是教材中演示自感现象的两个电路图，$L_{1}$ 和 $L_{2}$ 为电感线圈。实验时，断开开关 $S_{1}$ 瞬间，灯 $A_{1}$ 突然闪亮，随后逐渐变暗；闭合开关 $S_{2}$，灯 $A_{2}$ 逐渐变亮，而另一个相同的灯 $A_{3}$ 立即变亮，最终 $A_{2}$ 与 $A_{3}$ 的亮度相同。下列说法正确的是\xzanswer{C}

\twopicture[numstyle=arabic, labelA=2017北京理综19a, labelB=2017北京理综19b, wA=6.5cm, wB=6.5cm]
{figs/04a.png}
{figs/04b.png}

\fourchoices[answer=C]
{图 1 中，$A_{1}$ 与 $L_{1}$ 的电阻值相同}
{图 1 中，闭合 $S_{1}$，电路稳定后，$A_{1}$ 中电流大于 $L_{1}$ 中电流}
{图 2 中，变阻器 $R$ 与 $L_{2}$ 的电阻值相同}
{图 2 中，闭合 $S_{2}$ 瞬间，$L_{2}$ 中电流与变阻器 $R$ 中电流相等}

%% answer: C
%% memo:
\memoanswer{
断开开关 $S_{1}$ 瞬间，灯 $A_{1}$ 突然闪亮，由于线圈 $L_{1}$ 的自感，通过 $L_{1}$ 的电流逐渐减小，且通过 $A_{1}$，即自感电流会大于原来通过 $A_{1}$ 的电流，说明闭合 $S_{1}$，电路稳定时，通过 $A_{1}$ 的电流小于通过 $L_{1}$ 的电流，$L_{1}$ 的电阻小于 $A_{1}$ 的电阻，AB 错误；

闭合 $S_{2}$，电路稳定时，$A_{2}$ 与 $A_{3}$ 的亮度相同，说明两支路的电流相同，因此变阻器 $R$ 与 $L_{2}$ 的电阻值相同，C 正确；

闭合开关 $S_{2}$，$A_{2}$ 逐渐变亮，而 $A_{3}$ 立即变亮，说明 $L_{2}$ 中电流与变阻器 $R$ 中电流不相等，D 错误。
}

% ---------------------------------------------------------------------------
% 第 5 题：2011 年北京卷第 19 题（上游 试卷/2011/北京理综/北京理综.tex 中 number 为 19 的一题）
% 配图：05.png←试卷/2011/北京理综/figs/19.png
\item
%% number: 19
%% paperName: 2011年北京理综第 19 题 6 分
%% typeId: 1
%% type: 单选题
%% chapter: 电磁感应
%% point: 自感与互感，涡流
%% method: 无
%% score: 6
%% degree: 450
%% duplicateId: 0
%% body:
\source{2011·北京卷·19}某同学为了验证断电自感现象，自己找来带铁心的线圈 $L$、小灯泡 A、开关 S 和电池组 $E$，用导线将它们连接成如图所示的电路。检查电路后，闭合开关 S，小灯泡发光；再断开开关 S，小灯泡仅有不显著的延时熄灭现象。虽经多次重复，仍未见老师演示时出现的小灯泡闪亮现象，他冥思苦想找不出原因。你认为最有可能造成小灯泡未闪亮的原因是\xzanswer{C}

\onepicture[w=6cm]{figs/05.png}

\fourchoices[answer=C]
{电源的内阻较大}
{小灯泡电阻偏大}
{线圈电阻偏大}
{线圈的自感系数较大}
%% answer: C
%% memo:
\memoanswer{
小灯泡没有出现闪亮现象是因为断电后电路中的小灯泡两端电压太小。因断电后电路与电源脱离关系，线圈与灯泡组成闭合回路，故电源内阻大小对自感无影响，故 A 错误；若小灯泡电阻偏大，则分得的电压就大，这有助于出现闪亮现象，故 B 错误；若线圈电阻偏大，在自感电动势一定的情况下，线圈内阻上的电压偏大，相应灯泡两端的电压就偏小，这不利于出现闪亮现象，故 C 正确；线圈自感系数越大，产生的自感电动势越大，这有利于闪亮现象的出现，故 D 错误。
}

% ---------------------------------------------------------------------------
% 第 6 题：2008 年江苏卷第 8 题（上游 试卷/2008/江苏/江苏.tex 中 number 为 8 的一题）
% 配图：06.png←试卷/2008/江苏/figs/08.png
\item
%% number: 8
%% paperName: 2008年江苏高考第 8 题 4 分
%% typeId: 2
%% type: 多选题
%% chapter: 电磁感应
%% point: 自感与互感，涡流
%% method: 无
%% score: 4
%% degree: 650
%% duplicateId: 0
%% body:
\source{2008·江苏卷·8}如图所示的电路中，三个相同的灯泡 a、b、c 和电感 $L_{1}$、$L_{2}$ 与直流电源连接，电感的电阻忽略不计，电键 K 从闭合状态突然断开时，下列判断正确的有\xzanswer{AD}

\onepicture[w=5.5cm]{figs/06.png}

\fourchoices[answer=AD]
{a 先变亮，然后逐渐变暗}
{b 先变亮，然后逐渐变暗}
{c 先变亮，然后逐渐变暗}
{b、c 都逐渐变暗}

%% answer: AD
%% memo:
\memoanswer{
设电键 K 闭合、电路稳定时每个灯泡中的电流为 $I$。电感的电阻不计，三个灯泡两端的电压都等于电源电动势，此时 $L_{1}$ 中的电流等于 b、c 中电流之和，为 $2I$，$L_{2}$ 中的电流为 $I$。

K 断开瞬间，电源被断开，$L_{1}$、$L_{2}$ 中的电流不能突变，二者与三个灯泡组成回路：a 中的电流为 $2I$（方向与原来相反），大于原来的电流，故 a 先变亮，然后逐渐变暗，A 正确；b 中的电流为 $2I - I = I$，c 中的电流为 $I$，都不大于原来的电流，故 b、c 都不会变亮，而是逐渐变暗，B、C 错误，D 正确。

故选 AD。
}

% ---------------------------------------------------------------------------
% 第 7 题：2010 年北京卷第 19 题（上游 试卷/2010/北京理综/北京理综.tex 中 number 为 19 的一题）
% 配图：07.png←试卷/2010/北京理综/figs/19.png；07a.png←试卷/2010/北京理综/figs/19a.png；07b.png←试卷/2010/北京理综/figs/19b.png；07c.png←试卷/2010/北京理综/figs/19c.png；07d.png←试卷/2010/北京理综/figs/19d.png
\item
%% number: 19
%% paperName: 2010年北京理综第 19 题 6 分
%% typeId: 1
%% type: 单选题
%% chapter: 电磁感应
%% point: 自感与互感，涡流
%% method: 无
%% score: 6
%% degree: 450
%% duplicateId: 0
%% body:
\source{2010·北京卷·19}在如图所示的电路中，两个相同的小灯泡 $\mathrm{L}_{1}$ 和 $\mathrm{L}_{2}$ 分别串联一个带铁芯的电感线圈 $L$ 和一个滑动变阻器 $R$。闭合开关 $S$ 后，调整 $R$，使 $\mathrm{L}_{1}$ 和 $\mathrm{L}_{2}$ 发光的亮度一样，此时流过两个灯泡的电流均为 $I$。然后，断开 $S$。若 $t'$ 时刻再闭合 $S$，则在 $t'$ 前后的一小段时间内，正确反映流过 $\mathrm{L}_{1}$ 的电流 $i_{1}$、流过 $\mathrm{L}_{2}$ 的电流 $i_{2}$ 随时间 $t$ 变化的图像是\xzanswer{B}
\onepicture[align=right, w=5.5cm]{figs/07.png}
\fourchoices[answer=B, ispicture=true, h=2.2cm]
{figs/07a.png}
{figs/07b.png}
{figs/07c.png}
{figs/07d.png}
%% answer: B
%% memo:
\memoanswer{
$t'$ 时刻闭合 S 时，$\mathrm{L}_{1}$ 与带铁芯的电感线圈 $L$ 串联，线圈产生的自感电动势阻碍电流增大，$i_{1}$ 从 0 逐渐增大到 $I$，A 错误，B 正确；

$\mathrm{L}_{2}$ 与滑动变阻器 $R$ 串联，该支路没有自感作用，$i_{2}$ 立即增大到 $I$，方向不变，C、D 错误。

故选 B。
}

% ---------------------------------------------------------------------------
% 第 8 题：2010 年江苏卷第 4 题（上游 试卷/2010/江苏/江苏.tex 中 number 为 4 的一题）
% 配图：08a.png←试卷/2010/江苏/figs/04a.png；08b.png←试卷/2010/江苏/figs/04b.png；08c.png←试卷/2010/江苏/figs/04c.png；08d.png←试卷/2010/江苏/figs/04d.png；08e.png←试卷/2010/江苏/figs/04e.png
\item
%% number: 4
%% paperName: 2010年江苏高考第 4 题 3 分
%% typeId: 1
%% type: 单选题
%% chapter: 交流电
%% point: LC振荡回路
%% method: 无
%% score: 3
%% degree: 650
%% duplicateId: 0
%% body:
\source{2010·江苏卷·4}如图所示的电路中，电源的电动势为 $E$，内阻为 $r$，电感 $L$ 的电阻不计，电阻 $R$ 的阻值大于灯泡 D 的阻值，在 $t=0$ 时刻闭合开关 S，经过一段时间后，在 $t=t_{1}$ 时刻断开 S，下列表示 A、B 两点间电压 $U_{AB}$ 随时间 $t$ 变化的图象中，正确的是\xzanswer{B}

\onepicture[w=4.5cm]{figs/08a.png}

\fourchoices[answer=B, ispicture=true, h=2.4cm]
{figs/08b.png}
{figs/08c.png}
{figs/08d.png}
{figs/08e.png}

%% answer: B
%% memo:
\memoanswer{
开关闭合时，线圈由于自感对电流的阻碍作用，可看做电阻，线圈电阻逐渐减小，并联电路电阻逐渐减小，电压 $U_{AB}$ 逐渐减小；开关闭合后再断开时，线圈的感应电流与原电流方向相同，形成回路，灯泡的电流与原电流方向相反，并逐渐减小到 0，所以正确选项 B。
}

% ---------------------------------------------------------------------------
% 第 9 题：2025 年北京卷第 9 题（上游 试卷/2025/北京/北京.tex 中 number 为 9 的一题）
% 配图：09.png←试卷/2025/北京/figs/09.png
\item
%% number: 9
%% paperName: 2025年北京高考第 9 题 3 分
%% typeId: 1
%% type: 单选题
%% chapter: 电磁感应
%% point: 自感与互感，涡流
%% method: 无
%% score: 3
%% degree: 450
%% duplicateId: 0
%% body:
\source{2025·北京卷·9}如图所示，线圈自感系数为 $L$，电容器电容为 $C$，电源电动势为 $E$，A$_{1}$、A$_{2}$ 和 A$_{3}$ 是三个相同的小灯泡。开始时，开关 S 处于断开状态。忽略线圈电阻和电源内阻，将开关 S 闭合，下列说法正确的是\xzanswer{C}

\onepicture[w=4.00cm, h=2.85cm]{figs/09.png}

\fourchoices[answer=C]
{闭合瞬间，A$_{1}$ 与 A$_{3}$ 同时亮起}
{闭合后，A$_{2}$ 亮起后亮度不变}
{稳定后，A$_{1}$ 与 A$_{3}$ 亮度一样}
{稳定后，电容器的电荷量是 $CE$}

%% answer: C
%% memo:
\memoanswer{
A．闭合开关瞬间，电容器 $C$ 相当于通路，线圈 $L$ 相当于断路，所以 A$_{1}$、A$_{2}$ 瞬间亮起，A$_{3}$ 逐渐变亮，A 错误；

B．闭合开关后，电容器充电，充电完成后相当于断路，所以 A$_{2}$ 亮一下后熄灭，B 错误；

C．稳定后，电容器相当于断路，线圈相当于短路，所以 A$_{1}$、A$_{3}$ 串联，所以一样亮，C 正确；

D．稳定后，电容器与 A$_{3}$ 并联，两端电压等于 A$_{3}$ 两端电压，由于线圈电阻和电源内阻忽略不计，且 A$_{1}$、A$_{3}$ 串联，A$_{3}$ 两端电压为 $\frac{1}{2}E$，根据 $Q = CU$，可得电容器的电荷量等于 $\frac{1}{2}CE$，D 错误。

故选 C。
}

% ---------------------------------------------------------------------------
% 第 10 题：2025 年浙江1月选考第 9 题（上游 试卷/2025/1月浙江/1月浙江.tex 中 number 为 9 的一题）
% 配图：10a.png←试卷/2025/1月浙江/figs/09a.png；10b.png←试卷/2025/1月浙江/figs/09b.png；10c.png←试卷/2025/1月浙江/figs/09c.png；10d.png←试卷/2025/1月浙江/figs/09d.png
\item
%% number: 9
%% paperName: 2025年1月浙江高考第 9 题 3 分
%% typeId: 1
%% type: 单选题
%% chapter: 电磁感应
%% point: 自感与互感，涡流
%% method: 无
%% score: 3
%% degree: 450
%% duplicateId: 0
%% body:
\source{2025·浙江1月选考·9}新能源汽车日趋普及，其能量回收系统可将制动时的动能回收再利用，当制动过程中回收系统的输出电压（$U$）比动力电池所需充电电压（$U_{0}$）低时，不能直接充入其中。在下列电路中，通过不断打开和闭合开关 S，实现由低压向高压充电，其中正确的是\xzanswer{B}

\fourchoices[answer=B, ispicture=true, h=3.00cm]
{figs/10a.png}
{figs/10b.png}
{figs/10c.png}
{figs/10d.png}

%% answer: B
%% memo:
\memoanswer{
A．该电路中当开关 S 断开时，整个电路均断开，则不能给电池充电，选项 A 错误；

B．该电路中当 S 闭合时稳定时，线圈 L 中有电流通过，当 S 断开时 L 产生自感电动势阻碍电流减小，L 相当电源，电源 U 与 L 中的自感电动势共同加在电池两端，且此时二极管能导通，从而实现给高压充电，选项 B 正确；

C．该电路中当 S 闭合时稳定时，线圈 L 中有电流通过，但当 S 断开时 L 也与电路断开，还是只有回收系统的电压 U 加在充电电池两端，则不能实现给高压充电，选项 C 错误；

D．该电路中当 S 闭合时稳定时，线圈 L 中有电流通过，但当 S 断开时电源 U 也断开，只有 L 产生的自感电动势相当电源加在充电电池两端，则不能实现给高压充电，选项 D 错误。

故选 B。
}

\end{enumerate}

\section{涡流、电磁阻尼和电磁驱动}
\begin{enumerate}[resume]
% ---------------------------------------------------------------------------
% 第 11 题：2024 年湖北卷第 1 题（上游 试卷/2024/湖北/湖北.tex 中 number 为 1 的一题）
% 配图：无
\item
%% number: 1
%% paperName: 2024年湖北高考第 1 题 4 分
%% typeId: 1
%% type: 单选题
%% chapter: 电磁感应
%% point: 自感与互感，涡流
%% method: 无
%% score: 4
%% degree: 950
%% duplicateId: 0
%% body:
\source{2024·湖北卷·1}《梦溪笔谈》中记录了一次罕见的雷击事件：房屋被雷击后，屋内的银饰、宝刀等金属熔化了，但是漆器、刀鞘等非金属却完好（原文为：有一木格，其中杂贮诸器，其漆器银扣者，银悉熔流在地，漆器曾不焦灼。有一宝刀，极坚钢，就刀室中熔为汁，而室亦俨然）。导致金属熔化而非金属完好的原因可能为 \xzanswer{C} 

\fourchoices[answer=C]
{摩擦}
{声波}
{涡流}
{光照}

%% answer: C
%% memo:
\memoanswer{
在雷击事件中金属和非金属都经历了摩擦，声波和光照的影响，而金属能够因电磁感应产生涡流非金属不能，因此可能原因为涡流。

故选 C。
}

% ---------------------------------------------------------------------------
% 第 12 题：2020 年全国卷Ⅱ第 14 题（上游 试卷/2020/全国理综Ⅱ/全国理综Ⅱ.tex 中 number 为 1 的一题）
% 配图：12.png←试卷/2020/全国理综Ⅱ/figs/01.png
\item
%% number: 1
%% paperName: 2020年全国理综Ⅱ第 1 题 6 分
%% typeId: 1
%% type: 单选题
%% chapter: 电磁感应
%% point: 电磁感应现象
%% method: 无
%% score: 6
%% degree: 950
%% duplicateId: 0
%% body:
\source{2020·全国卷Ⅱ·14}管道高频焊机可以对由钢板卷成的圆管的接缝实施焊接。焊机的原理如图所示，圆管通过一个接有高频交流电源的线圈，线圈所产生的交变磁场使圆管中产生交变电流，电流产生的热量使接缝处的材料熔化将其焊接。焊接过程中所利用的电磁学规律的发现者为\xzanswer{D}

\onepicture[w=6.50cm]{figs/12.png}

\fourchoices[answer=D]
{库仑}
{霍尔}
{洛伦兹}
{法拉第}
%% answer: D
%% memo:
\memoanswer{
由题意可知，圆管为金属导体，导体内部自成闭合回路，且有电阻，当周围的线圈中产生出交变磁场时，就会在导体内部感应出涡电流，电流通过电阻要发热。该过程利用的原理是电磁感应现象，其发现者为法拉第。故选 D。
}

% ---------------------------------------------------------------------------
% 第 13 题：2024 年甘肃卷第 6 题（上游 试卷/2024/甘肃/甘肃.tex 中 number 为 6 的一题）
% 配图：13.png←试卷/2024/甘肃/figs/06.png
\item
%% number: 6
%% paperName: 2024年甘肃高考第 6 题 4 分
%% typeId: 1
%% type: 单选题
%% chapter: 电磁感应
%% point: 自感与互感，涡流
%% method: 无
%% score: 4
%% degree: 850
%% duplicateId: 0
%% body:
\source{2024·甘肃卷·6}工业上常利用感应电炉冶炼合金，装置如图所示。当线圈中通有交变电流时，下列说法正确的是\xzanswer{B}

\onepicture[w=3.50cm]{figs/13.png}

\fourchoices[answer=B]
{金属中产生恒定感应电流}
{金属中产生交变感应电流}
{若线圈匝数增加，则金属中感应电流减小}
{若线圈匝数增加，则金属中感应电流不变}

%% answer: B
%% memo:
\memoanswer{
AB．当线圈中通有交变电流时，感应电炉金属内的磁通量也不断随之变化，金属中产生交变感应电流，A 错误，B 正确；

CD．若线圈匝数增加，根据电磁感应定律可知，感应电动势增大，则金属中感应电流变大，CD 错误。

故选 B。
}

% ---------------------------------------------------------------------------
% 第 14 题：2014 年江苏卷第 7 题（上游 试卷/2014/江苏/江苏.tex 中 number 为 7 的一题）
% 配图：14.png←试卷/2014/江苏/figs/07.png
\item
%% number: 7
%% paperName: 2014年江苏高考第 7 题 4 分
%% typeId: 2
%% type: 多选题
%% chapter: 交流电
%% point: 交流电
%% method: 无
%% score: 4
%% degree: 450
%% duplicateId: 0
%% body:
\source{2014·江苏卷·7}如图所示，在线圈上端放置一盛有冷水的金属杯，现接通交流电源，过了几分钟，杯内的水沸腾起来。若要缩短上述加热时间，下列措施可行的有\xzanswer{AB}

\onepicture[align=right, w=4.50cm]{figs/14.png}

\fourchoices[answer=AB]
{增加线圈的匝数}
{提高交流电源的频率}
{将金属杯换为瓷杯}
{取走线圈中的铁芯}

%% answer: AB
%% memo:
\memoanswer{
金属杯处在线圈产生的交变磁场中，杯壁内产生涡流，涡流的热效应使水温升高。要缩短加热时间，应增大涡流的热功率，即增大金属杯中的感应电动势。

A．增加线圈的匝数，线圈产生的磁场增强，穿过金属杯的磁通量变化率增大，涡流增大，A 正确；

B．提高交流电源的频率，磁场变化加快，磁通量变化率增大，涡流增大，B 正确；

C．瓷杯是绝缘体，不能产生涡流，无法加热，C 错误；

D．取走铁芯，磁场大为减弱，涡流减小，D 错误。

故选 AB。
}

% ---------------------------------------------------------------------------
% 第 15 题：2022 年浙江1月选考第 13 题（上游 试卷/2022/1月浙江/1月浙江.tex 中 number 为 13 的一题）
% 配图：15.png←试卷/2022/1月浙江/figs/13.png
\item
%% number: 13
%% paperName: 2022年1月浙江选考第 13 题 3 分
%% typeId: 1
%% type: 单选题
%% chapter: 电磁感应
%% point: 法拉第电磁感应定律
%% method: 无
%% score: 3
%% degree: 450
%% duplicateId: 0
%% body:
\source{2022·浙江1月选考·13}如图所示，将一通电螺线管竖直放置，螺线管内部形成方向竖直向上、磁感应强度大小 $B=kt$ 的匀强磁场，在内部用绝缘轻绳悬挂一与螺线管共轴的金属薄圆管，其电阻率为 $\rho$、高度为 $h$、半径为 $r$、厚度为 $d$（$d\ll r$），则 \xzanswer{C}

\onepicture[h=4.00cm, align=right]{figs/15.png}

\fourchoices[answer=C]
{从上向下看，圆管中的感应电流为逆时针方向}
{圆管的感应电动势大小为 $\frac{k\pi r^{2}}{h}$}
{圆管的热功率大小为 $\frac{\pi dhk^{2}r^{3}}{2\rho}$}
{轻绳对圆管的拉力随时间减小}

%% answer: C
%% memo:
\memoanswer{
	A．穿过圆管的磁通量向上逐渐增加，则根据楞次定律可知，从上向下看，圆管中的感应电流为顺时针方向，选项 A 错误；

	B．圆管的感应电动势大小为

	$E=\frac{\Delta B}{\Delta t}\pi r^{2}=k\pi r^{2}$

	选项 B 错误；

	C．圆管的电阻

	$R=\rho\frac{2\pi r}{dh}$

	圆管的热功率大小为

	$P=\frac{E^{2}}{R}=\frac{\pi dhk^{2}r^{3}}{2\rho}$

	选项 C 正确；

	D．根据左手定则可知，圆管中各段所受的安培力方向指向圆管的轴线，则轻绳对圆管的拉力的合力始终等于圆管的重力，不随时间变化，选项 D 错误。

	故选 C。
}

% ---------------------------------------------------------------------------
% 第 16 题：2025 年河南卷第 5 题（上游 试卷/2025/河南/河南.tex 中 number 为 5 的一题）
% 配图：16a.png←试卷/2025/河南/figs/05a.png；16b.png←试卷/2025/河南/figs/05b.png；16c.png←试卷/2025/河南/figs/05c.png；16d.png←试卷/2025/河南/figs/05d.png；16e.png←试卷/2025/河南/figs/05e.png
\item
%% number: 5
%% paperName: 2025年河南高考第 5 题 4 分
%% typeId: 1
%% type: 单选题
%% chapter: 电磁感应
%% point: 自感与互感，涡流
%% method: 无
%% score: 4
%% degree: 650
%% duplicateId: 0
%% body:
\source{2025·河南卷·5}如图，一金属薄片在力 $F$ 作用下自左向右从两磁极之间通过。当金属薄片中心运动到 N 极的正下方时，沿 N 极到 S 极的方向看，下列图中能够正确描述金属薄片内涡电流绕行方向的是\xzanswer{C}

\onepicture[w=3.60cm]{figs/16a.png}

\fourchoices[answer=C, ispicture=true, h=1.40cm]
{figs/16b.png}
{figs/16c.png}
{figs/16d.png}
{figs/16e.png}

%% answer: C
%% memo:
\memoanswer{
根据题意当金属薄片中心运动到 N 极正下方时，薄片右侧的磁通量在减小，左侧磁通量在增加，由于两极间的磁场竖直向下，根据楞次定律可知此时薄片右侧的涡电流方向为顺时针，薄片左侧的涡电流方向为逆时针。

故选 C。
}

% ---------------------------------------------------------------------------
% 第 17 题：2015 年山东卷第 17 题（上游 试卷/2015/山东理综/山东理综.tex 中 number 为 17 的一题）
% 配图：17.png←试卷/2015/山东理综/figs/17.png
\item
%% number: 17
%% paperName: 2015年山东理综第 17 题 6 分
%% typeId: 2
%% type: 多选题
%% chapter: 电磁感应
%% point: 右手定则
%% method: 无
%% score: 6
%% degree: 450
%% duplicateId: 0
%% body:
\source{2015·山东卷·17}如图，一均匀金属圆盘绕通过其圆心且与盘面垂直的轴逆时针匀速转动。现施加一垂直穿过圆盘的有界匀强磁场，圆盘开始减速。在圆盘减速过程中，以下说法正确的是\xzanswer{ABD}

\onepicture[h=2.6cm]{figs/17.png}

\fourchoices[answer=ABD]
{处于磁场中的圆盘部分，靠近圆心处电势高}
{所加磁场越强越易使圆盘停止转动}
{若所加磁场反向，圆盘将加速转动}
{若所加磁场穿过整个圆盘，圆盘将匀速转动}

%% answer: ABD
%% memo:
\memoanswer{
	由右手定则可知，处于磁场中的圆盘部分，靠近圆心处电势高，选项 A 正确；根据 $E=BLv$ 可知所加磁场越强，则感应电动势越大，感应电流越大，产生的电功率越大，消耗的机械能越快，则圆盘越容易停止转动，选项 B 正确；若加反向磁场，根据楞次定律可知安培力阻碍圆盘的转动，故圆盘仍减速转动，选项 C 错误；若所加磁场穿过整个圆盘则圆盘中无感应电流，不消耗机械能，圆盘匀速转动。选项 D 正确，选项 ABD 正确。
}

% ---------------------------------------------------------------------------
% 第 18 题：2015 年全国卷Ⅰ第 19 题（上游 试卷/2015/全国课标Ⅰ/全国课标Ⅰ.tex 中 number 为 19 的一题）
% 配图：18.svg←试卷/2015/全国课标Ⅰ/figs/19.svg
\item
%% number: 19
%% paperName: 2015年全国课标Ⅰ第 19 题 6 分
%% typeId: 2
%% type: 多选题
%% chapter: 电磁感应
%% point: 楞次定律推广
%% method: 无
%% score: 6
%% degree: 450
%% duplicateId: 0
%% body:
\source{2015·全国卷Ⅰ·19}1824 年，法国科学家阿拉果完成了著名的“圆盘实验”。实验中将一铜圆盘水平放置，在其中心正上方用柔软细线悬挂一枚可以自由旋转的磁针，如图所示。实验中发现，当圆盘在磁针的磁场中绕过圆盘中心的竖直轴旋转时，磁针也随着一起转动起来，但略有滞后。下列说法正确的是\xzanswer{AB}

\onepicture[h=4.3cm]{\includesvg{figs/18}}

\fourchoices[answer=AB]
{圆盘上产生了感应电动势}
{圆盘内的涡电流产生的磁场导致磁针转动}
{在圆盘转动的过程中，磁针的磁场穿过整个圆盘的磁通量发生了变化}
{圆盘中的自由电子随圆盘一起运动形成电流，此电流产生的磁场导致磁针转动}

%% answer: AB
%% memo:
\memoanswer{
A．圆盘可看作由无数根沿半径方向的金属条组成，圆盘转动时，金属条在磁针的磁场中做切割磁感线运动，圆盘上产生感应电动势，A 正确；

B．圆盘中的感应电动势在盘内形成涡电流，涡电流的磁场与磁针相互作用，使磁针跟着圆盘转动起来（电磁驱动），B 正确；

C．磁针的磁感线从 N 极出发回到 S 极，在圆盘所在平面上向下穿过与向上穿过的部分相互抵消，穿过整个圆盘的磁通量为零，圆盘转动过程中保持不变，C 错误；

D．圆盘整体呈电中性，转动时自由电子与带正电的原子实一起运动，不会形成电流，D 错误。

故选 AB。
}

% ---------------------------------------------------------------------------
% 第 19 题：2014 年广东卷第 15 题（上游 试卷/2014/广东理综/广东理综.tex 中 number 为 15 的一题）
% 配图：19.png←试卷/2014/广东理综/figs/15.png
\item
%% number: 15
%% paperName: 2014年广东理综第 15 题 4 分
%% typeId: 1
%% type: 单选题
%% chapter: 电磁感应
%% point: 楞次定律推广
%% method: 无
%% score: 4
%% degree: 650
%% duplicateId: 0
%% body:
\source{2014·广东卷·15}如图所示，上下开口、内壁光滑的铜管 P 和塑料管 Q 竖直放置，小磁块先后在两管中从相同高度处由静止释放，并落至底部，则小磁块\xzanswer{C}
\onepicture[align=right, w=3.2cm]{figs/19.png}
\fourchoices[answer=C]
{在 P 和 Q 中都做自由落体运动}
{在两个下落过程中的机械能都守恒}
{在 P 中的下落时间比在 Q 中的长}
{落至底部时在 P 中的速度比在 Q 中的大}
%% answer: C
%% memo:
\memoanswer{
铜管是导体，小磁块在铜管 P 中下落时，管壁中产生感应电流（涡流），感应电流的磁场阻碍小磁块与铜管的相对运动，小磁块受到向上的磁场力，加速度小于 $g$，不做自由落体运动，且部分机械能转化为电能（最终转化为内能），机械能不守恒；塑料管不导电，不产生感应电流，小磁块在 Q 中做自由落体运动，机械能守恒。故 A、B 错误。

小磁块在 P 中下落的加速度小于在 Q 中的加速度，从相同高度由静止下落，在 P 中的下落时间比在 Q 中的长，落至底部时在 P 中的速度比在 Q 中的小，C 正确，D 错误。

故选 C。
}

% ---------------------------------------------------------------------------
% 第 20 题：2014 年全国大纲卷第 20 题（上游 试卷/2014/全国理综大纲/全国理综大纲.tex 中 number 为 20 的一题）
% 配图：无
\item
%% number: 20
%% paperName: 2014年全国理综大纲第 20 题 6 分
%% typeId: 1
%% type: 单选题
%% chapter: 电磁感应
%% point: 楞次定律推广
%% method: 无
%% score: 6
%% degree: 650
%% duplicateId: 0
%% body:
\source{2014·全国大纲卷·20}很多相同的绝缘铜圆环沿竖直方向叠放，形成一很长的竖直圆筒。一条形磁铁沿圆筒的中心轴竖直放置，其下端与圆筒上端开口平齐。让条形磁铁从静止开始下落。条形磁铁在圆筒中的运动速率\xzanswer{C}
\fourchoices[answer=C]
{均匀增大}
{先增大，后减小}
{逐渐增大，趋于不变}
{先增大，再减小，最后不变}
%% answer: C
%% memo:
\memoanswer{
小磁铁在圆环中下落的过程由于穿过圆环的磁通量发生变化，所以在圆环中会产生感应电流，阻碍磁铁的下落，又因为圆筒很长，所以当磁铁受到的阻力等于重力时，其速度保持趋于不变，选项 C 正确。
}

% ---------------------------------------------------------------------------
% 第 21 题：2023 年全国乙卷第 17 题（上游 试卷/2023/全国乙卷/全国乙卷.tex 中 number 为 4 的一题）
% 配图：21a.png←试卷/2023/全国乙卷/figs/04a.png；21b.png←试卷/2023/全国乙卷/figs/04b.png
\item
%% number: 4
%% paperName: 2023年全国高考乙卷第 4 题 6 分
%% typeId: 1
%% type: 单选题
%% chapter: 电磁感应
%% point: 电磁感应中图象
%% method: 无
%% score: 6
%% degree: 550
%% duplicateId: 0
%% body:
\source{2023·全国乙卷·17}一学生小组在探究电磁感应现象时，进行了如下比较实验。用图（a）所示的缠绕方式，将漆包线分别绕在几何尺寸相同的有机玻璃管和金属铝管上，漆包线的两端与电流传感器接通。两管皆竖直放置，将一很小的强磁体分别从管的上端由静止释放，在管内下落至管的下端。实验中电流传感器测得的两管上流过漆包线的电流 $I$ 随时间 $t$ 的变化分别如图（b）和图（c）所示，分析可知\xzanswer{A}
\twopicture[showlabel=false, labelA=2023全国乙卷04a, labelB=2023全国乙卷04b, hA=4.2cm, wB=9.6cm]
{figs/21a.png}
{figs/21b.png}
\fourchoices[answer=A]
{图（c）是用玻璃管获得的图像}
{在铝管中下落，小磁体做匀变速运动}
{在玻璃管中下落，小磁体受到的电磁阻力始终保持不变}
{用铝管时测得的电流第一个峰到最后一个峰的时间间隔比用玻璃管时的短}
%% answer: A
%% memo:
\memoanswer{
A．强磁体在铝管中运动，铝管会形成涡流，玻璃是绝缘体故强磁体在玻璃管中运动，玻璃管不会形成涡流。强磁体在铝管中加速后很快达到平衡状态，做匀速直线运动，而玻璃管中的磁体则一直做加速运动，故由图像可知图（c）的脉冲电流峰值不断增大，说明强磁体的速度在增大，与玻璃管中磁体的运动情况相符，A 正确；

B．在铝管中下落，脉冲电流的峰值一样，磁通量的变化率相同，故小磁体做匀速运动，B 错误；

C．在玻璃管中下落，玻璃管为绝缘体，线圈的脉冲电流峰值增大，电流不断在变化，故小磁体受到的电磁阻力在不断变化，C 错误；

D．强磁体分别从管的上端由静止释放，在铝管中，磁体在线圈间做匀速运动，玻璃管中磁体在线圈间做加速运动，故用铝管时测得的电流第一个峰到最后一个峰的时间间隔比用玻璃管时的长，D 错误。

故选 A。
}

% ---------------------------------------------------------------------------
% 第 22 题：2017 年全国卷Ⅰ第 18 题（上游 试卷/2017/全国理综Ⅰ/全国理综Ⅰ.tex 中 number 为 18 的一题）
% 配图：22.svg←试卷/2017/全国理综Ⅰ/figs/18.svg；22a.png←试卷/2017/全国理综Ⅰ/figs/18a.png；22b.png←试卷/2017/全国理综Ⅰ/figs/18b.png；22c.png←试卷/2017/全国理综Ⅰ/figs/18c.png；22d.png←试卷/2017/全国理综Ⅰ/figs/18d.png
\item
%% number: 18
%% paperName: 2017年全国理综Ⅰ第 18 题 6 分
%% typeId: 1
%% type: 单选题
%% chapter: 电磁感应
%% point: 楞次定律推广
%% method: 无
%% score: 6
%% degree: 650
%% duplicateId: 0
%% body:
\source{2017·全国卷Ⅰ·18}扫描隧道显微镜（STM）可用来探测样品表面原子尺寸上的形貌，为了有效隔离外界振动对 STM 的扰动，在圆底盘周边沿其径向对称地安装若干对紫铜薄板，并施加磁场来快速衰减其微小振动，如图所示，无扰动时，按下列四种方案对紫铜薄板施加恒磁场；出现扰动后，对于紫铜薄板上下及其左右振动的衰减最有效的方案是\xzanswer{A}

\onepicture[h=3.8cm]{\includesvg{figs/22}}

\fourchoices[answer=A, ispicture=true, h=1.8cm]
{figs/22a.png}
{figs/22b.png}
{figs/22c.png}
{figs/22d.png}

%% answer: A
%% memo:
\memoanswer{
	施加磁场来衰减振动是利用电磁阻尼实现的，即振动过程中发生电磁感应现象，产生的感应电流在磁场中受安培力阻碍其振动。A 图中紫铜薄板无论上下振动还是左右振动都会出现磁通量的变化，即产生电磁感应现象；B 图中紫铜薄板上下振动时没有磁通量的变化，不会产生电磁感应现象；C 图中紫铜薄板无论上下振动还是左右振动均不会产生电磁感应现象；D 图中紫铜薄板上下振动时不会产生电磁感应现象，故最有效的方案是 A，故 A 正确。
}

% ---------------------------------------------------------------------------
% 第 23 题：2021 年北京卷第 11 题（上游 试卷/2021/北京/北京.tex 中 number 为 11 的一题）
% 配图：23.png←试卷/2021/北京/figs/11.png
\item
%% number: 11
%% paperName: 2021年北京高考第 11 题 3 分
%% typeId: 1
%% type: 单选题
%% chapter: 电磁感应
%% point: 楞次定律推广
%% method: 无
%% score: 3
%% degree: 650
%% duplicateId: 0
%% body:
\source{2021·北京卷·11}某同学搬运如图所示的磁电式电流表时，发现表针晃动剧烈且不易停止。按照老师建议，该同学在两接线柱间接一根导线后再次搬运，发现表针晃动明显减弱且能很快停止。下列说法正确的是\xzanswer{D}

\onepicture[w=4cm]{figs/23.png}

\fourchoices[answer=D]
{未接导线时，表针晃动过程中表内线圈不产生感应电动势}
{未接导线时，表针晃动剧烈是因为表内线圈受到安培力的作用}
{接上导线后，表针晃动过程中表内线圈不产生感应电动势}
{接上导线后，表针晃动减弱是因为表内线圈受到安培力的作用}

%% answer: D
%% memo:
\memoanswer{
A．未接导线时，表针晃动过程中导线切割磁感线，表内线圈会产生感应电动势，故 A 错误；

B．未接导线时，未连成闭合回路，没有感应电流，所以不受安培力，故 B 错误；

CD．接上导线后，表针晃动过程中表内线圈产生感应电动势，根据楞次定律可知，表针晃动减弱是因为表内线圈受到安培力的作用，故 C 错误 D 正确。

故选 D。
}

% ---------------------------------------------------------------------------
% 第 24 题：2025 年北京卷第 10 题（上游 试卷/2025/北京/北京.tex 中 number 为 10 的一题）
% 配图：24.png←试卷/2025/北京/figs/10.png
\item
%% number: 10
%% paperName: 2025年北京高考第 10 题 3 分
%% typeId: 1
%% type: 单选题
%% chapter: 电磁感应
%% point: 楞次定律推广
%% method: 无
%% score: 3
%% degree: 450
%% duplicateId: 0
%% body:
\source{2025·北京卷·10}绝缘的轻质弹簧上端固定，下端悬挂一个磁铁。将磁铁从弹簧原长位置由静止释放，磁铁开始振动，由于空气阻力的影响，振动最终停止。现将一个闭合铜线圈固定在磁铁正下方的桌面上（如图所示），仍将磁铁从弹簧原长位置由静止释放，振动最终也停止。则\xzanswer{D}

\onepicture[w=2.10cm, h=4.40cm]{figs/24.png}

\fourchoices[answer=D]
{有无线圈，磁铁经过相同的时间停止运动}
{磁铁靠近线圈时，线圈有扩张趋势}
{磁铁离线圈最近时，线圈受到的安培力最大}
{有无线圈，磁铁和弹簧组成的系统损失的机械能相同}

%% answer: D
%% memo:
\memoanswer{
A．有线圈时，磁铁受到电磁阻尼的作用，振动更快停止，故 A 错误；

B．根据楞次定律，磁铁靠近线圈时，线圈的磁通量增大，此时线圈有缩小的趋势，故 B 错误；

C．磁铁离线圈最近时，此时磁铁与线圈的相对速度为零，感应电动势为零，感应电流为零，线圈受到的安培力为零，故 C 错误；

D．分析可知有无线圈时，根据平衡条件最后磁铁静止后弹簧的伸长量相同，由于磁铁和弹簧组成的系统损失的机械能为磁铁减小的重力势能减去此时弹簧的弹性势能，故系统损失的机械能相同，故 D 正确。

故选 D。
}

\end{enumerate}

\section{实验题与计算题}
\begin{enumerate}[resume]
% ---------------------------------------------------------------------------
% 第 25 题：2015 年江苏卷第 11 题（上游 试卷/2015/江苏理综/江苏理综.tex 中 number 为 11 的一题）
% 配图：25a.png←试卷/2015/江苏理综/figs/11a.png；25b.png←试卷/2015/江苏理综/figs/11b.png
\item
%% number: 11
%% paperName: 2015年江苏理综第 11 题 10 分
%% typeId: 4
%% type: 实验题
%% chapter: 电磁感应
%% point: 电磁感应中的变加速
%% method: 无
%% score: 10
%% degree: 450
%% duplicateId: 0
%% body:
\source{2015·江苏卷·11}某同学探究小磁铁在铜管中下落时受电磁阻尼作用的运动规律。实验装置如图所示，打点计时器的电源为 $50\UHz$ 的交流电。

\onepicture[h=4.4cm, align=right]{figs/25a.png}

（1）下列实验操作中，不正确的有\tkanswer[2.0cm]{CD}。

（A）将铜管竖直地固定在限位孔的正下方

（B）纸带穿过限位孔，压在复写纸下面

（C）用手捏紧磁铁保持静止，然后轻轻地松开让磁铁下落

（D）在磁铁下落的同时接通打点计时器的电源

（2）该同学按正确的步骤进行实验（记为“实验①”），将磁铁从管口处释放，打出一条纸带，取开始下落的一段，确定一合适的点为 O 点，每隔一个计时点取一个计数点，标为 1，2，$\ldots$，8。用刻度尺量出各计数点的相邻两计时点到 O 点的距离，记录在纸带上，如图所示。

\onepicture[h=1.6cm]{figs/25b.png}

计算相邻计时点间的平均速度 $\bar v$，粗略地表示各计数点的速度，抄入下表。请将表中的数据补充完整。

\begin{center}
\begin{tabular}{|c|c|c|c|c|c|c|c|c|}
\hline
位置 & 1 & 2 & 3 & 4 & 5 & 6 & 7 & 8 \\
\hline
$\bar v$（$\Ucms$） & 24.5 & 33.8 & 37.8 & \tkanswer[1.2cm]{39.0} & 39.5 & 39.8 & 39.8 & 39.8 \\
\hline
\end{tabular}
\end{center}

（3）分析上表的实验数据可知：在这段纸带记录的时间内，磁铁运动速度的变化情况是\tkanswer[4.0cm]{逐渐增大到 $39.8\Ucms$}；磁铁受到阻尼作用的变化情况是\tkanswer[4.0cm]{逐渐增大到等于重力}。

（4）该同学将装置中的铜管更换为相同尺寸的塑料管，重复上述实验操作（记为“实验②”），结果表明磁铁下落的运动规律与自由落体运动规律几乎相同。请问实验②是为了说明什么？对比实验①和②的结果可得到什么结论？

\tkanswer[10cm]{为了说明磁铁在塑料管中几乎不受阻尼作用。磁铁在铜管中受到的阻尼作用主要是电磁阻尼作用。}

%% answer:
\jdanswer{
（1）CD；
（2）39.0；
（3）逐渐增大到 $39.8\Ucms$，逐渐增大到等于重力；
（4）为了说明磁铁在塑料管中几乎不受阻尼作用。磁铁在铜管中受到的阻尼作用主要是电磁阻尼作用。
}
%% memo:
\memoanswer{
（1）打点计时器使用时应先接通电源，待打点稳定后再释放纸带，D 错误；释放前应手持纸带上端，使纸带竖直、磁铁静止在管口，再松开纸带，而不应用手捏住磁铁，C 错误；A、B 的操作正确。故不正确的是 CD。

（2）打点周期 $T = 0.02\Us$，计数点 4 两侧相邻两计时点到 O 点的距离分别为 $4.04\Ucm$ 和 $5.60\Ucm$，则
\[ \bar v_{4} = \frac{(5.60 - 4.04)\Ucm}{2 \times 0.02\Us} = 39.0\Ucms \]

（3）由表中数据可知，磁铁运动的速度逐渐增大，增大得越来越慢，最后达到 $39.8\Ucms$ 后基本不变；即磁铁的加速度逐渐减小到零，由牛顿第二定律可知，磁铁受到的阻尼作用逐渐增大，最后等于重力。

（4）塑料管不导电，磁铁在其中下落几乎做自由落体运动，说明除电磁阻尼外，磁铁受到的其他阻力（如摩擦、空气阻力）可以忽略；对比实验①和②可知，磁铁在铜管中受到的阻尼作用主要是电磁阻尼作用。
}

% ---------------------------------------------------------------------------
% 第 26 题：2021 年北京卷第 19 题（上游 试卷/2021/北京/北京.tex 中 number 为 19 的一题）
% 配图：26a.png←试卷/2021/北京/figs/19a.png；26b.png←试卷/2021/北京/figs/19b.png；26_an.png←试卷/2021/北京/figs/19_an.png
\item
%% number: 19
%% paperName: 2021年北京高考第 19 题 10 分
%% typeId: 6
%% type: 计算题
%% chapter: 其他
%% point: 物理思想
%% method: 无
%% score: 10
%% degree: 550
%% duplicateId: 0
%% body:
\source{2021·北京卷·19}类比是研究问题的常用方法。

\begin{enumerate}
	\item
	情境 1：物体从静止开始下落，除受到重力作用外，还受到一个与运动方向相反的空气阻力 $f = kv$（$k$ 为常量）的作用。其速率 $v$ 随时间 $t$ 的变化规律可用方程 $G - kv = m\frac{\Delta v}{\Delta t}$（①式）描述，其中 $m$ 为物体质量，$G$ 为其重力。求物体下落的最大速率 $v_{\mathrm{m}}$。
	\item
	情境 2：如图 \ref{2021北京19a} 所示，电源电动势为 $E$，线圈自感系数为 $L$，电路中的总电阻为 $R$。闭合开关 S，发现电路中电流 $I$ 随时间 $t$ 的变化规律与情境 1 中物体速率 $v$ 随时间 $t$ 的变化规律类似。类比①式，写出电流 $I$ 随时间 $t$ 变化的方程；并在图 \ref{2021北京19b} 中定性画出 $I-t$ 图线。

	\twopicture[numA=1, numB=2, labelA=2021北京19a, labelB=2021北京19b, wA=4.6cm, wB=3.8cm]
	{figs/26a.png}{figs/26b.png}

	\item
	类比情境 1 和情境 2 中的能量转化情况，完成下表。

	\begin{center}
	\begin{tblr}{|c|c|}
	\hline
	情境 1 & 情境 2 \\
	\hline
	物体重力势能的减少量 & \\
	\hline
	物体动能的增加量 & \\
	\hline
	 & 电阻 $R$ 上消耗的电能 \\
	\hline
	\end{tblr}
	\end{center}
\end{enumerate}

%% answer: 
\jdanswer{
\begin{enumerate}
	\item
	$v_{\mathrm{m}} = \frac{G}{k}$

	\item
	$E - RI = L\frac{\Delta I}{\Delta t}$

	\item
	电源提供的电能，线圈磁场能的增加量，克服阻力做功消耗的机械能

\end{enumerate}
}
%% memo:
\memoanswer{
（1）当物体下落速度达到最大速度 $v_{\mathrm{m}}$ 时，加速度为零，有
$G = kv_{\mathrm{m}}$
得 $v_{\mathrm{m}} = \frac{G}{k}$

（2）类比①式，由闭合电路的欧姆定律有
$E - RI = L\frac{\Delta I}{\Delta t}$

闭合开关后，线圈产生的自感电动势阻碍电流增大，电流从零开始逐渐增大，且增大得越来越慢，电路稳定后自感现象消失，电流趋于 $\frac{E}{R}$，$I-t$ 图线如下

\onepicture[w=4.5cm]{figs/26_an.png}

（3）两情境中能量转化的对应关系如下表

\begin{center}
\begin{tblr}{|c|c|}
\hline
情境 1 & 情境 2 \\
\hline
 & 电源提供的电能 \\
\hline
 & 线圈磁场能的增加量 \\
\hline
克服阻力做功消耗的机械能 & \\
\hline
\end{tblr}
\end{center}
}

% ---------------------------------------------------------------------------
% 第 27 题：2022 年重庆卷第 13 题（上游 试卷/2022/重庆/重庆.tex 中 number 为 13 的一题）
% 配图：27.png←试卷/2022/重庆/figs/13.png
\item
%% number: 13
%% paperName: 2022年重庆高考第 13 题 12 分
%% typeId: 6
%% type: 计算题
%% chapter: 电磁感应
%% point: 法拉第电磁感应定律
%% method: 无
%% score: 12
%% degree: 650
%% duplicateId: 0
%% body:
\source{2022·重庆卷·13}某同学以金属戒指为研究对象，探究金属物品在变化磁场中的热效应。如图所示，戒指可视为周长为 $L$、横截面积为 $S$、电阻率为 $\rho$ 的单匝圆形线圈，放置在匀强磁场中，磁感应强度方向垂直于戒指平面。若磁感应强度大小在 $\Delta t$ 时间内从 0 均匀增加到 $B_{0}$，求：
\begin{enumerate}
	\item
	戒指中的感应电动势和电流；
	\item
	戒指中电流的热功率。
\end{enumerate}

\onepicture[w=4.00cm, align=right]{figs/27.png}

%% answer:
\jdanswer{
\begin{enumerate}
	\item
	$E = \frac{B_{0}L^{2}}{4\pi \Delta t}$，$I = \frac{B_{0}LS}{4\pi \rho \Delta t}$

	\item
	$P = \frac{B_{0}^{2}L^{3}S}{16\pi^{2}\rho(\Delta t)^{2}}$

\end{enumerate}
}
%% memo:
\memoanswer{
（1）设戒指的半径为 $r$，则有

$L = 2\pi r$

磁感应强度大小在 $\Delta t$ 时间内从 0 均匀增加到 $B_{0}$，产生的感应电动势为

$E = \frac{B_{0}}{\Delta t} \cdot \pi r^{2}$

可得 $E = \frac{B_{0}L^{2}}{4\pi \Delta t}$

戒指的电阻为 $R = \frac{\rho L}{S}$，则戒指中的感应电流为 $I = \frac{E}{R} = \frac{B_{0}LS}{4\pi \rho \Delta t}$

（2）戒指中电流的热功率为

$P = I^{2}R = \frac{B_{0}^{2}L^{3}S}{16\pi^{2}\rho(\Delta t)^{2}}$
}

% ---------------------------------------------------------------------------
% 第 28 题：2016 年天津卷第 12 题（上游 试卷/2016/天津理综/天津理综.tex 中 number 为 12 的一题）
% 配图：28.png←试卷/2016/天津理综/figs/12.png
\item
%% number: 12
%% paperName: 2016年天津理综第 12 题 20 分
%% typeId: 6
%% type: 计算题
%% chapter: 电磁感应
%% point: 电磁感应中的变加速
%% method: 无
%% score: 20
%% degree: 350
%% duplicateId: 0
%% body:
\source{2016·天津卷·12}电磁缓速器是应用于车辆上以提高运行安全性的辅助制动装置，其工作原理是利用电磁阻尼作用减缓车辆的速度。电磁阻尼作用可以借助如下模型讨论：如图所示，将形状相同的两根平行且足够长的铝条固定在光滑斜面上，斜面与水平方向夹角为 $\theta$。一质量为 $m$ 的条形磁铁滑入两铝条间，恰好匀速穿过，穿过时磁铁两端面与两铝条的间距始终保持恒定，其引起电磁感应的效果与磁铁不动、铝条相对磁铁运动相同。磁铁端面是边长为 $d$ 的正方形，由于磁铁距离铝条很近，磁铁端面正对两铝条区域的磁场均可视为匀强磁场，磁感应强度为 $B$，铝条的高度大于 $d$，电阻率为 $\rho$。为研究问题方便，铝条中只考虑与磁铁正对部分的电阻和磁场，其他部分电阻和磁场可忽略不计，假设磁铁进入铝条间以后，减少的机械能完全转化为铝条的内能，重力加速度为 $g$。
\begin{enumerate}
	\item
	求铝条中与磁铁正对部分的电流 $I$；
	\item
	若两铝条的宽度均为 $b$，推导磁铁匀速穿过铝条间时速度 $v$ 的表达式；
	\item
	在其他条件不变的情况下，仅将两铝条更换为宽度 $b'>b$ 的铝条，磁铁仍以速度 $v$ 进入铝条间，试简要分析说明磁铁在铝条间运动时的加速度和速度如何变化。
\end{enumerate}
\onepicture[h=6.2cm, align=right]{figs/28.png}

%% answer: 
\jdanswer{
\begin{enumerate}
	\item
	$I=\frac{mg\sin \theta }{2Bd}$

	\item
	$v=\frac{\rho mg\sin \theta }{2B^{2}d^{2}b}$

	\item
	磁铁做加速度逐渐减小的减速运动，直到 $F'=mg\sin\theta$ 时，磁铁重新达到平衡状态，将再次以较小的速度匀速下滑。

\end{enumerate}
}
%% memo:
\memoanswer{
（1）磁铁在铝条间运动时，两根铝条受到的安培力大小相等均为 $F_{\text{安}}$，有 $F_{\text{安}} = IdB$ ①

磁铁受到沿斜面向上的作用力为 $F$，其大小有 $F = 2F_{\text{安}}$ ②

磁铁匀速运动时受力平衡，则有 $F - mg \sin \theta = 0$ ③

联立①②③式可得 $I = \frac{mg \sin \theta}{2Bd}$ ④

（2）磁铁穿过铝条时，在铝条中产生的感应电动势为 $E$，有 $E = Bdv$ ⑤

铝条与磁铁正对部分的电阻为 $R$，由电阻定律有 $R = \rho \frac{d}{db}$ ⑥

由欧姆定律有 $I = \frac{E}{R}$ ⑦

联立④⑤⑥⑦式可得 $v = \frac{\rho mg \sin \theta}{2B^{2}d^{2}b}$ ⑧

（3）磁铁以速度 $v$ 进入铝条间，恰好做匀速运动时，磁铁受到沿斜面向上的作用力 $F$，联立①②⑤⑥⑦式可得 $F = \frac{2B^{2}d^{2}bv}{\rho}$ ⑨

当铝条的宽度 $b' > b$ 时，磁铁以速度 $v$ 进入铝条间时，磁铁受到的作用力变为 $F'$，有 $F' = \frac{2B^{2}d^{2}b'v}{\rho}$

可见，$F' > F = mg\sin\theta$，磁铁所受到的合力方向沿斜面向上，获得与运动方向相反的加速度，磁铁将减速下滑，此时加速度最大。之后，随着运动速度减小，$F'$ 也随着减小，磁铁所受的合力也减小，由于磁铁加速度与所受到的合力成正比，磁铁的加速度逐渐减小。综上所述，磁铁做加速度逐渐减小的减速运动，直到 $F' = mg\sin\theta$ 时，磁铁重新达到平衡状态，将再次以较小的速度匀速下滑。
}

\end{enumerate}

\end{document}
